\documentclass[11pt,a4paper]{article}

% --- Codificación y fuentes ---
% xelatex maneja UTF-8 nativamente; no necesita inputenc/fontenc
\usepackage[spanish]{babel}

% --- Layout ---
\usepackage[margin=2.2cm,top=2.5cm,bottom=2.5cm]{geometry}
\usepackage{setspace}
\setstretch{1.15}
\usepackage{microtype}

% --- Color, gráficos, hipervínculos ---
\usepackage{xcolor}
\definecolor{darkblue}{RGB}{26,58,92}
\definecolor{midblue}{RGB}{42,90,140}
\definecolor{lightgray}{RGB}{245,245,245}
\definecolor{darkgray}{RGB}{80,80,80}
\definecolor{accentred}{RGB}{200,50,50}
\usepackage{graphicx}
\usepackage{float}
\usepackage[hidelinks,colorlinks=true,linkcolor=midblue,urlcolor=midblue,citecolor=midblue]{hyperref}

% --- Tablas ---
\usepackage{booktabs}
\usepackage{array}
\usepackage{tabularx}
\usepackage{multirow}
\renewcommand{\arraystretch}{1.15}

% --- Matemáticas y unidades ---
\usepackage{amsmath,amssymb}
\usepackage{siunitx}

% --- Pseudocódigo (estilo del ejemplo del usuario) ---
\usepackage{algorithm}
\usepackage{algpseudocode}
\algrenewcommand\algorithmicrequire{\textbf{Input:}}
\algrenewcommand\algorithmicensure{\textbf{Output:}}
\algrenewcommand\textproc{\textsc}
% Comando "pardo" para for paralelos (estilo PRAM)
\algnewcommand\algorithmicpardo{\textbf{\color{accentred}pardo}}
\algnewcommand\Pardo{\algorithmicpardo}
\algdef{SE}[FOR]{ParFor}{EndParFor}[1]%
  {\algorithmicfor\ #1\ \algorithmicpardo}%
  {\algorithmicend\ \algorithmicfor}

% --- Código fuente ---
\usepackage{listings}
\lstset{
  basicstyle=\ttfamily\footnotesize,
  backgroundcolor=\color{lightgray},
  frame=single,
  framerule=0pt,
  rulecolor=\color{lightgray},
  keywordstyle=\color{midblue}\bfseries,
  commentstyle=\color{darkgray}\itshape,
  stringstyle=\color{accentred},
  showstringspaces=false,
  breaklines=true,
  columns=fullflexible,
  language=Python,
  morekeywords={comm,bcast,Bcast,Scatterv,Gatherv,MPI,mpirun,Wtime,Barrier},
  xleftmargin=8pt,
  xrightmargin=4pt,
  aboveskip=8pt,
  belowskip=8pt
}

% --- Encabezados y pie ---
\usepackage{fancyhdr}
\pagestyle{fancy}
\fancyhf{}
\fancyhead[L]{\small\textcolor{darkgray}{Paralelización de KNN con MPI}}
\fancyhead[R]{\small\textcolor{darkgray}{Programación Paralela}}
\fancyfoot[C]{\thepage}
\renewcommand{\headrulewidth}{0.4pt}
\renewcommand{\footrulewidth}{0pt}

% --- Títulos de sección ---
\usepackage{titlesec}
\titleformat{\section}{\Large\bfseries\color{darkblue}}{\thesection}{0.6em}{}
\titleformat{\subsection}{\large\bfseries\color{midblue}}{\thesubsection}{0.5em}{}
\titleformat{\subsubsection}{\normalsize\bfseries\color{midblue}}{\thesubsubsection}{0.5em}{}
\titlespacing*{\section}{0pt}{18pt}{8pt}
\titlespacing*{\subsection}{0pt}{14pt}{6pt}

% --- Caja para destacar versiones beta ---
\usepackage[most]{tcolorbox}
\newtcolorbox{betabox}[1]{
  colback=lightgray,
  colframe=midblue,
  fonttitle=\bfseries\color{white},
  coltitle=white,
  title=#1,
  boxrule=0pt,
  arc=2pt,
  left=8pt,right=8pt,top=4pt,bottom=4pt,
  before skip=10pt, after skip=10pt
}

% ============================================================
\begin{document}
% ============================================================

% --- Portada ---
\begin{titlepage}
\centering
\vspace*{2cm}
{\Huge\bfseries\color{darkblue} Paralelización de KNN con MPI\par}
\vspace{0.4cm}
{\Large\itshape sobre el dataset \texttt{load\_digits}\par}
\vspace{2cm}
\rule{0.8\textwidth}{0.8pt}\\[0.5cm]
\begin{tabular}{rl}
\textbf{Curso:} & Programación Paralela \\[2pt]
\textbf{Fecha:} & Abril 2026 \\[2pt]
\textbf{Hardware:} & Apple M2 Pro (8 P-cores + 4 E-cores) \\[2pt]
\textbf{Software:} & Open MPI 5.0.9, Python, mpi4py \\[2pt]
\textbf{Repositorio:} & \texttt{knn-mpi-parallel/} \\
\end{tabular}\\[0.3cm]
\rule{0.8\textwidth}{0.8pt}\\
\vspace{1.5cm}
\begin{minipage}{0.85\textwidth}
\textbf{Resumen.} Se paraleliza el algoritmo KNN del archivo \texttt{knn\_digits\_sec.py} sobre el dataset \texttt{load\_digits} (64 features, 10 clases) usando MPI. La estrategia sigue el DAG del enunciado: \texttt{bcast} del entrenamiento, \texttt{scatter} del testeo, cómputo local independiente y \texttt{gather} de predicciones. El proyecto se desarrolló en \textbf{tres versiones beta} documentando el proceso iterativo, incluyendo el descubrimiento y resolución de un problema de \emph{oversubscription} de threads BLAS común al combinar MPI con NumPy. La implementación final alcanza \textbf{6.14$\times$ de speedup} con 8 procesos y mantiene 76\,\% de eficiencia, con accuracy preservada idéntica al secuencial en todas las configuraciones.
\end{minipage}
\vfill
\end{titlepage}

\tableofcontents
\newpage

% ============================================================
\section{Resultados principales}
% ============================================================

La Tabla~\ref{tab:resumen} sintetiza el desempeño de la implementación final (v3) sobre el caso $n_{total}=40\,000$, $k=3$. La accuracy se mantiene idéntica al baseline secuencial en todas las configuraciones, y el speedup escala consistentemente con el número de procesos hasta llegar al límite de 8 P-cores físicos.

\begin{table}[H]
\centering
\caption{Resultados principales sobre Apple M2 Pro ($n_{total}=40\,000$, $k=3$).}
\label{tab:resumen}
\small
\begin{tabular}{r r r r r r}
\toprule
$p$ & $T$ (ms) & Speedup & Eficiencia & GFLOPs/s & Accuracy \\
\midrule
1 & 2\,362 & 1.00$\times$ & 100\,\% & 20.9 & 1.0000 \\
2 & 1\,134 & 2.08$\times$ & 104\,\% & 43.6 & 1.0000 \\
4 &    629 & 3.76$\times$ & 94\,\% & 78.9 & 1.0000 \\
8 &    385 & 6.14$\times$ & 76\,\% & 130.0 & 1.0000 \\
\bottomrule
\end{tabular}
\end{table}

La fracción serial estimada por el modelo de Amdahl es $f=2.24\,\%$, lo que confirma a KNN como un problema \emph{embarrassingly parallel} en su núcleo. La cantidad óptima de procesos estimada es $p^* \approx 16$ sobre el M2 Pro, limitada en la práctica por la cantidad de cores físicos y el sobrecosto de coordinación MPI por proceso.

% ============================================================
\section{Modelo PRAM y código en Python con directivas MPI}\label{sec:a}
% ============================================================

\subsection*{Modelo PRAM}

El algoritmo se modela como \textbf{CREW-PRAM} (\emph{Concurrent Read, Exclusive Write}) por las siguientes razones:

\begin{itemize}
\item \textbf{Concurrent Read:} todos los procesos leen el mismo conjunto de entrenamiento $X_{tr}$, replicado en cada \emph{rank} mediante \texttt{comm.Bcast}. La lectura concurrente sobre la misma estructura no genera conflictos.
\item \textbf{Exclusive Write:} cada proceso calcula y escribe predicciones únicamente para su fracción $n_{te}/p$ del conjunto de testeo. Las posiciones de escritura del vector global de predicciones son disjuntas por construcción.
\end{itemize}

% ─── ALGORITMO 1: Modelo PRAM (en función de n y p) ─────────────
\begin{algorithm}[H]
\small
\caption{KNN paralelo en modelo PRAM (en función de $n$ y $p$)}\label{alg:knn-pram}
\begin{algorithmic}[1]
\Require $X_{tr}[1..n_{tr},\,1..d]$, $y_{tr}[1..n_{tr}]$, $X_{te}[1..n_{te},\,1..d]$, $k$, $p$
\Ensure $y_{pred}[1..n_{te}]$
\Statex
\State \textit{// Fase 1: Difusión del entrenamiento}
\ParFor{$1 \leq q \leq p$}
    \State $X_{tr}^{(q)},\ y_{tr}^{(q)} \gets \textsc{Broadcast}(X_{tr},\ y_{tr})$
\EndParFor
\State \Comment{Costo: $O(\log p \cdot (\alpha + n_{tr}\cdot d\cdot \beta))$}
\Statex
\State \textit{// Fase 2: Partición del testeo}
\ParFor{$1 \leq q \leq p$}
    \State $X_{te}^{(q)} \gets X_{te}\bigl[\,(q-1)\lceil n_{te}/p\rceil + 1\ :\ q\lceil n_{te}/p\rceil\,\bigr]$
\EndParFor
\State \Comment{Costo: $O(p\cdot \alpha + n_{te}\cdot d\cdot \beta)$}
\Statex
\State \textit{// Fase 3: Cómputo de distancias y $k$-vecinos (núcleo paralelo)}
\ParFor{$1 \leq q \leq p$}
    \ParFor{$1 \leq j \leq \lceil n_{te}/p\rceil$}
        \For{$1 \leq \ell \leq n_{tr}$}
            \State $D[j,\ell] \gets \sum_{m=1}^{d}\bigl(X_{te}^{(q)}[j,m] - X_{tr}[\ell,m]\bigr)^{2}$
        \EndFor
        \State $K \gets \textsc{ArgPartition}(D[j,\,1..n_{tr}],\ k)$
        \State $y_{pred}^{(q)}[j] \gets \textsc{MajorityVote}(y_{tr}[K])$
    \EndParFor
\EndParFor
\State \Comment{Costo: $O(n_{tr} \cdot n_{te} / p)$ operaciones aritméticas}
\Statex
\State \textit{// Fase 4: Recolección de predicciones}
\ParFor{$1 \leq q \leq p$}
    \State $\textsc{Gather}(y_{pred}^{(q)} \to y_{pred})$
\EndParFor
\State \Comment{Costo: $O(p\cdot \alpha + n_{te}\cdot \beta)$}
\Statex
\State \Return $y_{pred}$
\end{algorithmic}
\end{algorithm}

El Algoritmo~\ref{alg:knn-pram} describe el procedimiento en el \textbf{modelo abstracto PRAM}, parametrizado en $n_{tr}$, $n_{te}$, $d$ y $p$. Cada fase declara su costo asintótico en función de estas variables, lo que permite verificar directamente la complejidad total $O(\log p \cdot n_{tr} d\, \beta) + O(n_{tr} n_{te}/p) + O(p\,\alpha)$. La distancia euclidiana aparece expandida como $\sum_{m=1}^{d}(\cdot)^2$ para hacer visible el factor $d$ que reaparece en el conteo de FLOPs del inciso (d).

A continuación, el Algoritmo~\ref{alg:knn-mpi} muestra la \textbf{realización concreta en MPI}: la misma estructura lógica, pero escrita desde la perspectiva de un rank individual y usando las primitivas reales \texttt{Bcast}/\texttt{Scatterv}/\texttt{Gatherv} de la librería.

% ─── ALGORITMO 2: Implementación MPI ────────────────────────────
\begin{algorithm}[H]
\small
\caption{Implementación MPI del Algoritmo~\ref{alg:knn-pram} (vista por rank)}\label{alg:knn-mpi}
\begin{algorithmic}[1]
\Require $X_{tr} \in \mathbb{R}^{n_{tr}\times d}$, $y_{tr} \in \mathbb{Z}^{n_{tr}}$, $X_{te} \in \mathbb{R}^{n_{te}\times d}$, $k$, número de procesos $p$
\Ensure $y_{pred} \in \mathbb{Z}^{n_{te}}$ \Comment{etiquetas predichas}
\Statex
\State \textit{// Fase 1: Difusión del entrenamiento (Bcast)}
\ParFor{$i = 0$ \textbf{to} $p-1$}
    \State \textsc{Bcast}$(X_{tr}, y_{tr})$ desde rank $0$ a proceso $i$ \Comment{$O(\log p \cdot (\alpha + n_{tr} d \beta))$}
\EndParFor
\Statex
\State \textit{// Fase 2: Distribución del testeo (Scatter)}
\ParFor{$i = 0$ \textbf{to} $p-1$}
    \State $X_{te}^{(i)} \gets \textsc{Scatter}(X_{te},\ \text{bloque } i)$ \Comment{$O(p \cdot (\alpha + (n_{te}/p)\, d\, \beta))$}
\EndParFor
\Statex
\State \textit{// Fase 3: Cómputo local independiente (paralelo)}
\ParFor{$i = 0$ \textbf{to} $p-1$}
    \ParFor{$j = 1$ \textbf{to} $|X_{te}^{(i)}|$}
        \For{$\ell = 1$ \textbf{to} $n_{tr}$}
            \State $D[j,\ell] \gets \|X_{te}^{(i)}[j] - X_{tr}[\ell]\|^2$ \Comment{distancia al cuadrado}
        \EndFor
        \State $K \gets \textsc{ArgPartition}(D[j,:],\ k)[1{:}k]$ \Comment{índices de los $k$ menores}
        \State $L \gets y_{tr}[K]$
        \State $y_{pred}^{(i)}[j] \gets \textsc{MajorityVote}(L)$
    \EndParFor
\EndParFor
\State \Comment{Costo total de la fase 3: $O(n_{tr}\cdot n_{te}\,/\,p)$}
\Statex
\State \textit{// Fase 4: Recolección de predicciones (Gather)}
\ParFor{$i = 0$ \textbf{to} $p-1$}
    \State \textsc{Gather}$(y_{pred}^{(i)})$ hacia rank $0$ \Comment{$O(p \cdot (\alpha + k \beta))$}
\EndParFor
\Statex
\State \Return $y_{pred}$ \Comment{ensamblado en rank $0$}
\end{algorithmic}
\end{algorithm}

\subsection*{Mapeo del DAG a directivas MPI}

La Tabla~\ref{tab:dag-mpi} muestra la correspondencia uno-a-uno entre los nodos del DAG del enunciado, las directivas MPI usadas en el código y la complejidad teórica de cada operación, donde $\alpha$ es la latencia (s/mensaje) y $\beta$ el inverso del ancho de banda (s/byte).

\begin{table}[H]
\centering
\caption{Correspondencia entre el DAG, las directivas MPI y la complejidad asintótica.}\label{tab:dag-mpi}
\small
\begin{tabular}{l l l}
\toprule
\textbf{Nodo del DAG} & \textbf{Directiva MPI} & \textbf{Complejidad} \\
\midrule
Initial parameters & Carga en rank 0 & $O(n_{tr} d)$ \\
Broadcast superior & \texttt{comm.Bcast(X\_train, root=0)} & $O(\log p \cdot (\alpha + n_{tr} d \beta))$ \\
Scatter superior & \texttt{comm.Scatterv(X\_test, ...)} & $O(p \cdot (\alpha + (n_{te}/p)\, d\, \beta))$ \\
dist + sort + k-NN & \texttt{knn\_predict\_batch} (local) & $O(n_{tr}\cdot n_{te}\,/\,p)$ \\
Gather inferior & \texttt{comm.Gatherv(y\_pred\_local, ...)} & $O(p \cdot (\alpha + k \beta))$ \\
Apply majority & Ensamblado en rank 0 & $O(n_{te})$ \\
\bottomrule
\end{tabular}
\end{table}

Se utilizan \texttt{Scatterv} y \texttt{Gatherv} (versiones de longitud variable) en lugar de \texttt{Scatter}/\texttt{Gather} fijas para soportar $n_{te}$ no divisible por $p$ sin padding artificial. La diferencia de carga entre el rank más cargado y el menos cargado es como máximo de un punto de testeo.

\subsection*{Esqueleto del código}

\begin{lstlisting}[language=Python]
comm = MPI.COMM_WORLD
rank, size = comm.Get_rank(), comm.Get_size()

# 1. Rank 0 carga datos
if rank == 0:
    X_train, X_test, y_train, y_test = load_scaled_digits(args.n)

# 2. BCAST entrenamiento  -----  O(log p (alpha + n_tr beta))
comm.Bcast(X_train, root=0)
comm.Bcast(y_train, root=0)

# 3. SCATTERV testeo  ---------  O(p (alpha + n_te/p beta))
counts = [n_test // size + (1 if i < n_test % size else 0) for i in range(size)]
displs = [sum(counts[:i]) for i in range(size)]
comm.Scatterv([X_test, counts*d, displs*d, MPI.DOUBLE], X_test_local, root=0)

# 4. COMPUTO LOCAL  -----------  O(n_tr n_te / p)
y_pred_local = knn_predict_batch(X_test_local, X_train, y_train, k)

# 5. GATHERV predicciones  ----  O(p (alpha + k beta))
comm.Gatherv(y_pred_local, [y_pred, counts, displs, MPI.LONG], root=0)
\end{lstlisting}

% ============================================================
\section{Desarrollo en versiones beta}\label{sec:b}
% ============================================================

El proyecto se desarrolló iterativamente en tres versiones beta, cada una versionada en el repositorio. En las tres se valida que la \textbf{accuracy del modelo es idéntica al secuencial} para todos los valores de $p$, garantizando la corrección de la implementación.

% ── BETA v1 ──────────────────────────────────────────────────
\begin{betabox}{Beta v1 — Paralelización funcional con cómputo escalar}
\textbf{Archivo:} \texttt{src/knn\_mpi\_v1.py} \\[2pt]
\textbf{Objetivo:} validar la corrección del esquema MPI sin introducir variables.\\[2pt]
\textbf{Implementación:} sigue directamente el DAG con \texttt{Bcast}, \texttt{Scatterv} y \texttt{Gatherv}; el cómputo local mantiene los \emph{loops} Python escalares del código original (\texttt{euclidean\_distance} elemento a elemento, \texttt{argsort} completo, \texttt{Counter.most\_common}).
\end{betabox}

\textbf{Validación numérica} (dataset original, $n=1797$, mediana sobre 3 corridas):

\begin{table}[H]
\centering
\small
\begin{tabular}{r r r r}
\toprule
$p$ & $T$ (s) & Speedup & Accuracy \\
\midrule
seq & 1.2060 & 1.00$\times$ & 0.9833 \\
1 & 1.1994 & 1.01$\times$ & 0.9833 \\
2 & 0.6218 & 1.94$\times$ & 0.9833 \\
4 & 0.3245 & 3.72$\times$ & 0.9833 \\
8 & 0.1690 & 7.14$\times$ & 0.9833 \\
\bottomrule
\end{tabular}
\end{table}

Se obtiene \textbf{89\,\% de eficiencia a $p=8$}. El cómputo absorbe el 92--99\,\% del tiempo total, por lo que la comunicación es prácticamente despreciable. La accuracy de 0.9833 coincide con el secuencial, confirmando la implementación correcta de \texttt{bcast}/\texttt{scatter}/\texttt{gather}.

% ── BETA v2 ──────────────────────────────────────────────────
\begin{betabox}{Beta v2 — Vectorización del cómputo local}
\textbf{Archivo:} \texttt{src/knn\_mpi\_v2.py} \\[2pt]
\textbf{Objetivo:} eliminar el cuello de botella del Python interpretado en el cómputo de distancias.\\[2pt]
\textbf{Tres optimizaciones aplicadas:}
\begin{enumerate}
\item \textbf{Distancias por GEMM con identidad de Lagrange:}
\[
\|a-b\|^2 = \|a\|^2 + \|b\|^2 - 2\,a^\top b
\]
La matriz cruzada $X_{te}^{(i)} \cdot X_{tr}^\top$ se delega a BLAS. Se omite la $\sqrt{\cdot}$ porque preserva el orden y ahorra $n_{tr}\cdot n_{te}$ operaciones.
\item \textbf{Top-$k$ con \texttt{argpartition}:} $O(n)$ por fila, en vez de $O(n\log n)$ de \texttt{argsort}.
\item \textbf{Voto mayoritario vectorizado:} matriz \emph{one-hot} de votos + \texttt{argmax}, eliminando \texttt{Counter.most\_common}.
\end{enumerate}
\end{betabox}

\textbf{Resultado de la vectorización:} el tiempo de cómputo a $p=1$ cayó de 1.20\,s (v1) a \textbf{0.0058\,s} (v2) --- aproximadamente \textbf{207$\times$ más rápido}.

\textbf{Hallazgo importante.} Aunque la accuracy se mantuvo correcta, el speedup paralelo se degradó a partir de $p \geq 4$. A $n=40\,000$, $p=4$ daba $2.26\times$ y $p=8$ apenas $2.25\times$. El análisis de las repeticiones evidenció \textbf{varianza descontrolada} (CV hasta 100\,\%).

\textbf{Causa raíz identificada.} BLAS multi-threaded (Apple Accelerate en macOS) por defecto usa todos los cores disponibles. Al lanzar \texttt{mpirun -n 8} se generaban 8 procesos $\times$ 8 threads BLAS = \textbf{64 threads compitiendo por 8 cores físicos}. El \emph{scheduler} del sistema operativo introducía \emph{context-switching} constante, contaminando las mediciones.

% ── BETA v3 ──────────────────────────────────────────────────
\begin{betabox}{Beta v3 — Versión instrumentada para benchmarks}
\textbf{Archivo:} \texttt{src/knn\_mpi\_v3.py} \\[2pt]
\textbf{Objetivo:} producir mediciones reproducibles y limpias para el análisis del informe.\\[2pt]
\textbf{Cinco mejoras integradas:}
\begin{enumerate}
\item \textbf{Pinning de BLAS a 1 thread por proceso} mediante variables de entorno fijadas \emph{antes} de importar NumPy: \texttt{OMP\_NUM\_THREADS}, \texttt{MKL\_NUM\_THREADS}, \texttt{OPENBLAS\_NUM\_THREADS}, \texttt{VECLIB\_MAXIMUM\_THREADS}, \texttt{BLIS\_NUM\_THREADS}.
\item \texttt{comm.Barrier()} antes de cada fase para mediciones consistentes.
\item Repeticiones internas con \emph{warm-up} automático (la primera repetición se descarta).
\item Recolección de tiempos por rank (\texttt{t\_comp\_max/mean/min/std}) para detectar desbalance.
\item Salida directa a CSV con todas las métricas, lista para el análisis.
\end{enumerate}
\end{betabox}

\textbf{Validación del fix de BLAS:} con threads BLAS limitados a 1 por proceso, el speedup recuperó su comportamiento esperado (Tabla~\ref{tab:v2-vs-v3}).

\begin{table}[H]
\centering
\caption{Comparación del speedup a $p=8$ antes y después del fix de BLAS.}\label{tab:v2-vs-v3}
\small
\begin{tabular}{r r r r}
\toprule
$n$ & v2 (oversubscription) & v3 (con fix) & Mejora \\
\midrule
10\,000 & 4.05$\times$ & \textbf{4.98$\times$} & +23\,\% \\
40\,000 & 2.25$\times$ & \textbf{6.14$\times$} & +172\,\% \\
\bottomrule
\end{tabular}
\end{table}

Esta secuencia v1 $\to$ v2 $\to$ v3 documenta una buena práctica de programación paralela: \emph{al combinar MPI con bibliotecas que internamente usan threads (NumPy/BLAS, OpenMP, MKL), el threading interno debe limitarse a 1 para evitar oversubscription}. Es un error frecuente y mencionado explícitamente en la documentación oficial de mpi4py.

% ============================================================
\section{Mediciones de tiempo, validación teórica y $p$ óptimo}\label{sec:c}
% ============================================================

Las mediciones se obtuvieron con \texttt{experiments/run\_v3\_benchmark.sh} recorriendo el producto cartesiano $(n, p)$ con 5 repeticiones por configuración (15 para $n=40\,000$) y filtrado IQR de outliers. Se reportan \textbf{medianas} por su robustez frente a interrupciones del sistema operativo.

\subsection*{Tiempos de ejecución, cómputo y comunicación}

\begin{figure}[H]
\centering
\includegraphics[width=0.75\textwidth]{figures/time_vs_p.png}
\caption{Tiempo total en escala log-log para $n \in \{1797, 5000, 10000, 40000\}$. Las cuatro curvas muestran pendiente cercana a $-1$ hasta $p=4$ y se aplanan ligeramente a $p=8$. Cuanto mayor es $n$, mejor se mantiene la pendiente cerca de la ideal.}
\end{figure}

\begin{figure}[H]
\centering
\includegraphics[width=0.7\textwidth]{figures/time_breakdown.png}
\caption{Descomposición del tiempo total por fase para $n=40\,000$. El cómputo absorbe prácticamente todo el tiempo (verde); \texttt{bcast}/\texttt{scatter}/\texttt{gather} apenas son visibles. A $p=8$, la comunicación representa solo el 1.21\,\% del cómputo.}
\end{figure}

\subsection*{Tabla completa de tiempos}

\begin{table}[H]
\centering
\caption{Tiempos medianos en milisegundos, speedup, eficiencia y accuracy por configuración $(n,p)$.}
\small
\begin{tabular}{r r r r r r r r r}
\toprule
$n$ & $p$ & $T$ (ms) & $T_{\text{comp}}$ (ms) & $T_{\text{comm}}$ (ms) & C/Comp & Speedup & Ef. (\%) & Accuracy \\
\midrule
1\,797 & 1 & 4.47 & 4.45 & 0.02 & 0.4\,\% & 1.00 & 100.0 & 0.9833 \\
1\,797 & 2 & 2.60 & 2.53 & 0.07 & 2.8\,\% & 1.72 & 86.1 & 0.9833 \\
1\,797 & 4 & 1.53 & 1.39 & 0.13 & 9.4\,\% & 2.92 & 73.1 & 0.9833 \\
1\,797 & 8 & 1.09 & 0.85 & 0.24 & 28.2\,\% & 4.11 & 51.3 & 0.9833 \\
\addlinespace
5\,000 & 1 & 37.49 & 37.43 & 0.06 & 0.2\,\% & 1.00 & 100.0 & 0.9950 \\
5\,000 & 2 & 18.65 & 18.48 & 0.16 & 0.9\,\% & 2.01 & 100.5 & 0.9950 \\
5\,000 & 4 & 11.11 & 10.75 & 0.34 & 3.2\,\% & 3.37 & 84.4 & 0.9950 \\
5\,000 & 8 & 7.59 & 7.04 & 0.50 & 7.1\,\% & 4.94 & 61.7 & 0.9950 \\
\addlinespace
10\,000 & 1 & 121.00 & 120.85 & 0.10 & 0.1\,\% & 1.00 & 100.0 & 1.0000 \\
10\,000 & 2 & 61.21 & 60.94 & 0.17 & 0.3\,\% & 1.98 & 98.8 & 1.0000 \\
10\,000 & 4 & 36.06 & 35.42 & 0.51 & 1.4\,\% & 3.36 & 83.9 & 1.0000 \\
10\,000 & 8 & 24.30 & 23.31 & 0.79 & 3.4\,\% & 4.98 & 62.2 & 1.0000 \\
\addlinespace
40\,000 & 1 & 2362 & 2361 & 0.3 & 0.01\,\% & 1.00 & 100.0 & 1.0000 \\
40\,000 & 2 & 1134 & 1134 & 1.0 & 0.09\,\% & 2.08 & 104.2 & 1.0000 \\
40\,000 & 4 & 629 & 626 & 2.5 & 0.40\,\% & 3.76 & 94.0 & 1.0000 \\
40\,000 & 8 & 385 & 381 & 4.6 & 1.21\,\% & 6.14 & 76.5 & 1.0000 \\
\bottomrule
\end{tabular}
\end{table}

Hay dos observaciones notables. Primero, la \textbf{accuracy aumenta con $n$} (de 0.9833 a 1.0000) porque el dataset escalado contiene réplicas del original con ruido bajo, lo que hace al problema más fácil; lo importante es que sea \textbf{idéntica entre todos los $p$} para el mismo $n$, y eso se cumple. Segundo, a $n=40\,000$, $p=2$ muestra \textbf{eficiencia super-lineal} (104.2\,\%): efecto típico de \emph{cache}, el conjunto particionado entra mejor en L2/L3.

\subsection*{Validación teórica y normalización}

El modelo teórico que predice el tiempo total combina la región paralelizable, las comunicaciones colectivas y un sobrecosto por proceso:

\begin{equation}\label{eq:modelo}
T(p, n) \;=\; \underbrace{\frac{T_{\text{comp}}(n)}{p}}_{\text{paralelo}} \;+\; \underbrace{\alpha \log_2 p}_{\text{bcast}} \;+\; \underbrace{\gamma\,p}_{\text{coordinación}}
\end{equation}

Para hacer la expresión empíricamente ajustable se \textbf{normaliza fijando $T_{\text{comp}}(n) := T(1, n)$}, que es la medición secuencial real. Esto absorbe los costos absolutos de la máquina (FLOPs/s del CPU) en una sola constante medida y deja solo dos parámetros libres ($\alpha, \gamma$) por estimar de los datos.

El ajuste por mínimos cuadrados con restricciones físicas ($\alpha, \gamma \geq 0$) sobre $n=40\,000$ produce los valores de la Tabla~\ref{tab:ajuste}.

\begin{table}[H]
\centering
\caption{Parámetros ajustados del modelo teórico y modelo Amdahl alternativo.}\label{tab:ajuste}
\small
\begin{tabular}{l r l}
\toprule
\textbf{Parámetro} & \textbf{Valor estimado} & \textbf{Interpretación} \\
\midrule
$\alpha$ (latencia bcast) & $\approx 0$ ms & Despreciable en \emph{shared-memory} \\
$\gamma$ (sobrecosto/p)   & 9.13 ms/p     & Sincronización + gather \\
RMSE del ajuste           & 34 ms (3\,\%) & Buen ajuste \\
\addlinespace
$f$ (frac. serial Amdahl) & 2.24\,\%      & Algoritmo casi totalmente paralelo \\
$1/f$ (speedup asintótico)& 44.6$\times$  & Cota teórica máxima \\
\bottomrule
\end{tabular}
\end{table}

La predicción de Amdahl a $p=8$ es $S = 1/(f + (1-f)/8) = 6.91\times$, vs. $6.14\times$ medido (error $11\,\%$).

\begin{figure}[H]
\centering
\includegraphics[width=0.85\textwidth]{figures/theoretical_fit.png}
\caption{Comparación de los dos modelos contra los datos para $n=40\,000$. Ambos coinciden hasta $p\approx 10$. El modelo $\alpha\log p + \gamma p$ captura el sobrecosto creciente y predice un mínimo en $p^* \approx 16$, mientras Amdahl ignora ese término y por eso continúa decreciendo monotónicamente.}
\end{figure}

\subsection*{Cantidad óptima de procesos}

Derivando la Ecuación~\ref{eq:modelo} respecto a $p$ e igualando a cero:

\begin{equation}
\frac{\partial T}{\partial p} = -\frac{T_1}{p^2} + \frac{\alpha}{p\,\ln 2} + \gamma = 0
\end{equation}

Con $T_1 = 2362$ ms, $\alpha \approx 0$ y $\gamma = 9.13$ ms/p, el mínimo numérico se alcanza en:

\begin{equation}
\boxed{p^* \;\approx\; 16}
\end{equation}

con tiempo predicho $T(p^*) \approx 294$ ms y speedup esperado de $8.0\times$. Sin embargo, esta predicción asume cores físicos disponibles. \textbf{El M2 Pro solo tiene 8 P-cores}; más allá de $p=8$ se entra al territorio de los 4 E-cores ($\sim 40\,\%$ más lentos) o se vuelve a la oversubscription, donde la suposición lineal del modelo se rompe.

\textbf{Conclusión.} Dado el hardware del experimento, $p^* \in [8, 12]$ es la recomendación práctica. Los 8 P-cores son suficientes para alcanzar 76\,\% de eficiencia con un overhead manejable.

% ============================================================
\section{Speedup y FLOPs vs $p$}\label{sec:d}
% ============================================================

\begin{figure}[H]
\centering
\includegraphics[width=\textwidth]{figures/speedup_efficiency.png}
\caption{A la izquierda, speedup vs. $p$ para los cuatro tamaños del problema. La curva ideal $S=p$ aparece en línea discontinua. A la derecha, eficiencia paralela. Las curvas para $n$ grandes (5\,000, 10\,000, 40\,000) son muy similares hasta $p=4$; $n=1797$ cae más rápido por granularidad insuficiente. $n=40\,000$ es la única configuración que sostiene $>90\,\%$ de eficiencia hasta $p=4$.}
\end{figure}

\subsection*{Conteo de FLOPs según la fórmula del enunciado}

La distancia euclidiana entre dos vectores de dimensión $d$ se descompone como:

\begin{table}[H]
\centering
\small
\begin{tabular}{l r}
\toprule
\textbf{Operación} & \textbf{Cantidad} \\
\midrule
Restas $a_i - b_i$ & $d$ \\
Cuadrados $(a_i-b_i)^2$ & $d$ \\
Sumas acumulativas & $d - 1$ \\
Raíz cuadrada & $1$ \\
\midrule
\textbf{Total por par} & $\mathbf{3d + 1}$ \\
\bottomrule
\end{tabular}
\end{table}

Para \texttt{load\_digits} con $d=64$: \textbf{193 FLOPs por par de puntos}.

El total de FLOPs en la región paralelizable (entrenamiento) es:

\begin{equation}
\text{FLOPs}_{\text{total}} = (3d + 1) \cdot n_{tr} \cdot n_{te}
\end{equation}

Para $n=40\,000$ ($n_{tr}=32\,000$, $n_{te}=8\,000$): $\mathbf{49.4 \times 10^9}$ \textbf{FLOPs}. La métrica reportada como GFLOPs/s es $\text{FLOPs}_{\text{total}} / t_{comp}^{max}$ (rank más lento, cuello de botella).

\begin{figure}[H]
\centering
\includegraphics[width=0.75\textwidth]{figures/flops_vs_p.png}
\caption{Rendimiento de cómputo en GFLOPs/s. El crecimiento es prácticamente lineal con $p$ para $n$ grandes, lo que confirma la naturaleza embarrassingly parallel del cómputo. $n=10\,000$ alcanza el pico ($\approx 136$ GFLOPs/s a $p=8$) por mejor uso de \emph{cache} que $n=40\,000$.}
\end{figure}

\begin{table}[H]
\centering
\caption{Rendimiento de cómputo (GFLOPs/s) por configuración.}
\small
\begin{tabular}{r r r r r}
\toprule
$p$ & $n=1797$ & $n=5000$ & $n=10000$ & $n=40000$ \\
\midrule
1 & 22.4 & 20.6 & 25.6 & 20.9 \\
2 & 39.3 & 41.7 & 50.8 & 43.6 \\
4 & 70.9 & 71.6 & 87.7 & 78.9 \\
8 & 117.7 & 111.1 & 135.7 & 130.0 \\
\bottomrule
\end{tabular}
\end{table}

El crecimiento prácticamente lineal entre $p=1$ y $p=8$ (factor $\approx 6\times$ para $n=40\,000$) refleja directamente el speedup. La leve curvatura corresponde a la pérdida de eficiencia por el sobrecosto creciente identificado en la Sección~\ref{sec:c}.

% ============================================================
\section{Análisis de escalabilidad}\label{sec:e}
% ============================================================

\subsection*{Escalabilidad fuerte (\emph{strong scaling})}

Variación de $p$ manteniendo $n$ fijo. Las curvas de la figura de speedup muestran:

\begin{itemize}
\item \textbf{$n=1797$ (pequeño):} la eficiencia cae al 51\,\% en $p=8$. Razón: el tiempo de cómputo absoluto (4.5 ms) está cerca del piso de coordinación MPI, por lo que dividirlo más no compensa el overhead. \emph{No es paralelizable eficientemente más allá de $p=2$ para este tamaño.}
\item \textbf{$n=5000$ y $10\,000$ (medianos):} comportamiento prácticamente idéntico, eficiencia $\approx 84\,\%$ a $p=4$ y $\approx 62\,\%$ a $p=8$.
\item \textbf{$n=40\,000$ (grande):} la única que sostiene $>90\,\%$ a $p=4$ y aún 76\,\% a $p=8$. Es el régimen \emph{saludable}.
\end{itemize}

\subsection*{Escalabilidad débil (\emph{weak scaling})}

Mantener $n/p$ constante (carga constante por proceso). Aproximación a partir de los datos:

\begin{table}[H]
\centering
\small
\begin{tabular}{r r r r}
\toprule
$n$ & $p$ & $n/p$ & $T$ (ms) \\
\midrule
5\,000 & 1 & 5\,000 & 37.5 \\
10\,000 & 2 & 5\,000 & 61.2 \\
40\,000 & 8 & 5\,000 & 385 \\
\bottomrule
\end{tabular}
\end{table}

El tiempo \emph{no} se mantiene plano: crece sub-linealmente. Esto se debe a que el cómputo de KNN no escala como $n/p$ sino como $n_{tr}\cdot n_{te}/p$, y al replicar $n_{tr}$ completo en todos los procesos, doblar $n$ duplica el costo por proceso aunque también dupliquemos $p$. Por construcción, KNN tiene escalabilidad débil pobre cuando se mantiene la estructura ``todos contra todos''.

\subsection*{Análisis de isoeficiencia}

¿Cuánto debe crecer $n$ para mantener eficiencia constante al aumentar $p$? De la fórmula $E = 1/(1 + T_{\text{comm}} \cdot p / T_{\text{comp}})$:

\begin{table}[H]
\centering
\small
\begin{tabular}{r r r r}
\toprule
$p$ & $T_{\text{comm}}/T_{\text{comp}}$ & $E$ predicho & $E$ medido \\
\midrule
1 & 0.01\,\% & 1.000 & 1.000 \\
2 & 0.09\,\% & 0.998 & 1.042 \\
4 & 0.40\,\% & 0.984 & 0.940 \\
8 & 1.21\,\% & 0.912 & 0.765 \\
\bottomrule
\end{tabular}
\end{table}

La predicción basada solo en \emph{comm/comp} es \textbf{demasiado optimista} a $p=8$ (91\,\% vs. 76\,\% medido). La diferencia ($\sim 15\,\%$) viene de:

\begin{itemize}
\item Imbalance residual entre rank más rápido y más lento (\texttt{t\_comp\_std} no nulo).
\item \emph{Jitter} del sistema operativo (especialmente entre P-cores y E-cores en el M2 Pro).
\item \emph{Cache contention} al ejecutar múltiples procesos simultáneamente.
\end{itemize}

Para el algoritmo KNN como está formulado (con \texttt{bcast} del entrenamiento completo), la \textbf{isoeficiencia es aproximadamente cuadrática}: para mantener una eficiencia objetivo al duplicar $p$, $n$ debe crecer al menos como $n = \Theta(p \cdot \log p)$ debido al término de bcast, y en la práctica más cerca de $\Theta(p^2)$ una vez se considera que tanto $n_{tr}$ como $n_{te}$ crecen proporcionalmente a $n$.

\subsection*{Discusión de diseño: ¿por qué particionar test y no train?}\label{sec:test-vs-train}

Una pregunta natural es por qué la estrategia particiona el conjunto de testeo (más pequeño) en lugar del entrenamiento (más grande). La intuición ``más datos $=$ más trabajo paralelo'' es razonable, pero ignora un detalle crítico: \textbf{lo que se paraleliza no es el tamaño de los datos sino el número de tareas independientes}.

\subsubsection*{Origen de la independencia en KNN}

La estructura del algoritmo es:

\begin{lstlisting}[language=Python,basicstyle=\ttfamily\small]
para cada x_test en X_test:          # iteraciones independientes
    para cada x_train en X_train:    # debe ver TODOS los train
        d = distancia(x_test, x_train)
    k_nearest = top_k(distancias)
    prediccion = majority_vote(etiquetas[k_nearest])
\end{lstlisting}

Clasificar \texttt{x\_test[0]} es completamente independiente de clasificar \texttt{x\_test[1]}: el bucle exterior es \emph{embarrassingly parallel}. En cambio, el bucle interior es una \textbf{reducción}: para encontrar los $k$ vecinos de un solo punto de test, hay que examinar todos los puntos de train. Por tanto, el número natural de tareas independientes es exactamente $n_{te}$, no $n_{tr}$.

Duplicar $n_{tr}$ no genera tareas adicionales --- solo agranda el trabajo dentro de cada tarea existente.

\subsubsection*{Esquema alternativo: partition-train}

Existe el diseño alternativo, conocido como \emph{kNN distribuido data-parallel}:

\begin{itemize}
\item $X_{tr}$ se \texttt{Scatter} (cada rank tiene un subset disjunto de tamaño $n_{tr}/p$).
\item $X_{te}$ se \texttt{Bcast} (todos los ranks tienen el test completo).
\item Cada rank encuentra los $k$ \textbf{candidatos locales} para cada $x_{te}$ contra su subset de train.
\item Se hace un \texttt{Allreduce} con operador de \textbf{merge top-$k$} para combinar los $k$ globales por cada punto de test.
\end{itemize}

La Tabla~\ref{tab:test-vs-train} compara ambos esquemas en términos asintóticos.

\begin{table}[H]
\centering
\caption{Comparación de los dos esquemas de paralelización.}\label{tab:test-vs-train}
\small
\begin{tabular}{l l l}
\toprule
\textbf{Aspecto} & \textbf{Partition-test (este trabajo)} & \textbf{Partition-train (alternativo)} \\
\midrule
Memoria por rank & $O(n_{tr}\cdot d)$ train completo & $O(n_{tr}\cdot d / p)$ \\
Bcast inicial    & $O(\log p \cdot n_{tr} \cdot d \cdot \beta)$ del train & $O(\log p \cdot n_{te} \cdot d \cdot \beta)$ del test \\
Cómputo local    & $O(n_{tr}\cdot n_{te} / p)$ & $O(n_{tr}\cdot n_{te} / p)$ \\
Reducción final  & Gather etiquetas: $O(p \cdot k \cdot \beta)$ & \textbf{Allreduce top-$k$:} $O(p \cdot n_{te} \cdot k \cdot \beta)$ \\
Complejidad de código & Simple: \texttt{majority\_vote} local & Compleja: merge de listas-$k$ \\
\bottomrule
\end{tabular}
\end{table}

Hay tres observaciones clave de esta comparación:

\begin{enumerate}
\item \textbf{El cómputo es idéntico} en ambos esquemas. Esto es esperable: el costo es $n_{tr}\cdot n_{te}$ pares distancia $\div\,p$ procesadores, sin importar cómo se distribuyan los datos.
\item \textbf{Partition-train es más eficiente en memoria} (factor $p$). Solo importa cuando $X_{tr}$ no entra en RAM por nodo --- no es nuestro caso ($16$ MB cabe trivialmente en cualquier laptop moderna).
\item \textbf{Partition-train paga su ahorro de memoria en comunicación}: la reducción final transfiere $n_{te} \cdot k$ candidatos por rank, en vez de $n_{te}/p$ etiquetas finales. Para nuestros tamaños esto es $\sim n_{te} \cdot k$ vs. $\sim n_{te}/p$, un factor $p \cdot k$ peor.
\end{enumerate}

\subsubsection*{¿Cuándo conviene cada esquema?}

\begin{table}[H]
\centering
\small
\begin{tabular}{p{0.45\textwidth} p{0.45\textwidth}}
\toprule
\textbf{Partition-test (nuestro caso)} & \textbf{Partition-train} \\
\midrule
$X_{tr}$ cabe en RAM de cada nodo & $X_{tr}$ es enorme (no cabe) \\
Una máquina multicore o cluster pequeño & Cluster grande ($p \gg 10$) \\
Red rápida (shared-memory, InfiniBand) & Red lenta entre nodos \\
Implementación simple deseable & Se justifica complejidad extra \\
\bottomrule
\end{tabular}
\end{table}

\subsubsection*{Validación de la elección con datos del experimento}

La Tabla del análisis de tiempos muestra que a $n=40\,000$, $p=8$:

\begin{itemize}
\item Comunicación total: $4.6$ ms ($1.21\,\%$ del tiempo)
\item Bcast de train ($16$ MB): apenas un componente de esos $4.6$ ms
\end{itemize}

Es decir, replicar el train por Bcast \textbf{prácticamente no cuesta nada} en nuestro escenario. Cambiar a partition-train para ahorrar ese $1\,\%$ no compensaría la complejidad adicional del Allreduce top-$k$ ni el costo extra de comunicar $n_{te}\cdot k$ candidatos parciales. La elección de particionar test es la decisión correcta para el problema y hardware planteados.

\subsection*{Veredicto de escalabilidad}

KNN así implementado es:

\begin{itemize}
\item \textbf{Excelente para \emph{strong scaling}} cuando $n$ es grande ($\geq 10\,000$) y $p \leq$ número de cores físicos.
\item \textbf{Moderado para \emph{weak scaling}}: la replicación del entrenamiento limita la escalabilidad débil a problemas pequeños o medianos.
\item \textbf{Para escalar a clusters distribuidos} ($p \gg 10$), la mejor mejora algorítmica sería pasar a un esquema de \emph{partición del entrenamiento} (cada rank tiene un subset disjunto de $X_{tr}$) con \texttt{Allreduce} para combinar los $k$ mejores candidatos locales --- esto convierte el bcast en una operación distribuida y mejora la isoeficiencia.
\end{itemize}

% ============================================================
\section{Conclusiones}
% ============================================================

\begin{enumerate}
\item \textbf{Implementación validada.} La accuracy es idéntica al secuencial en las tres versiones beta para todas las configuraciones probadas.
\item \textbf{Desempeño.} Speedup de \textbf{6.14$\times$} con eficiencia de 76\,\% a $p=8$, $n=40\,000$, alcanzando \textbf{130 GFLOPs/s} sostenidos.
\item \textbf{Lección de ingeniería paralela.} Documentamos un caso real de \emph{oversubscription} de threads BLAS y su solución estándar (variables de entorno de \emph{pinning}), que resultó esencial para obtener mediciones reproducibles.
\item \textbf{Cantidad óptima de procesos.} $p^* \approx 16$ según el modelo teórico, pero limitado en la práctica a $p \approx 8$--$12$ por el hardware del experimento (8 P-cores).
\item \textbf{Escalabilidad.} El algoritmo es ideal para \emph{strong scaling} sobre datasets medianos-grandes en una sola máquina multicore. Para escalar a \emph{clusters} distribuidos requiere una variante con partición del entrenamiento.
\end{enumerate}

% ============================================================
\section*{Bibliografía}
\addcontentsline{toc}{section}{Bibliografía}
% ============================================================

\begin{enumerate}
\item Código secuencial base: \texttt{knn\_digits\_sec.py} (provisto por el curso).
\item Documentación oficial de mpi4py: \url{https://mpi4py.readthedocs.io/en/stable/}
\item Open MPI documentation: \url{https://www.open-mpi.org/doc/}
\item Repositorio del proyecto: \texttt{knn-mpi-parallel/} (incluye código, scripts de benchmark, notebook de análisis y CSV de resultados crudos).
\end{enumerate}

\end{document}
